% Options for packages loaded elsewhere
\PassOptionsToPackage{unicode}{hyperref}
\PassOptionsToPackage{hyphens}{url}
\documentclass[
]{article}
\usepackage{xcolor}
\usepackage[margin=1in]{geometry}
\usepackage{amsmath,amssymb}
\setcounter{secnumdepth}{-\maxdimen} % remove section numbering
\usepackage{iftex}
\ifPDFTeX
  \usepackage[T1]{fontenc}
  \usepackage[utf8]{inputenc}
  \usepackage{textcomp} % provide euro and other symbols
\else % if luatex or xetex
  \usepackage{unicode-math} % this also loads fontspec
  \defaultfontfeatures{Scale=MatchLowercase}
  \defaultfontfeatures[\rmfamily]{Ligatures=TeX,Scale=1}
\fi
\usepackage{lmodern}
\ifPDFTeX\else
  % xetex/luatex font selection
\fi
% Use upquote if available, for straight quotes in verbatim environments
\IfFileExists{upquote.sty}{\usepackage{upquote}}{}
\IfFileExists{microtype.sty}{% use microtype if available
  \usepackage[]{microtype}
  \UseMicrotypeSet[protrusion]{basicmath} % disable protrusion for tt fonts
}{}
\makeatletter
\@ifundefined{KOMAClassName}{% if non-KOMA class
  \IfFileExists{parskip.sty}{%
    \usepackage{parskip}
  }{% else
    \setlength{\parindent}{0pt}
    \setlength{\parskip}{6pt plus 2pt minus 1pt}}
}{% if KOMA class
  \KOMAoptions{parskip=half}}
\makeatother
\usepackage{color}
\usepackage{fancyvrb}
\newcommand{\VerbBar}{|}
\newcommand{\VERB}{\Verb[commandchars=\\\{\}]}
\DefineVerbatimEnvironment{Highlighting}{Verbatim}{commandchars=\\\{\}}
% Add ',fontsize=\small' for more characters per line
\usepackage{framed}
\definecolor{shadecolor}{RGB}{248,248,248}
\newenvironment{Shaded}{\begin{snugshade}}{\end{snugshade}}
\newcommand{\AlertTok}[1]{\textcolor[rgb]{0.94,0.16,0.16}{#1}}
\newcommand{\AnnotationTok}[1]{\textcolor[rgb]{0.56,0.35,0.01}{\textbf{\textit{#1}}}}
\newcommand{\AttributeTok}[1]{\textcolor[rgb]{0.13,0.29,0.53}{#1}}
\newcommand{\BaseNTok}[1]{\textcolor[rgb]{0.00,0.00,0.81}{#1}}
\newcommand{\BuiltInTok}[1]{#1}
\newcommand{\CharTok}[1]{\textcolor[rgb]{0.31,0.60,0.02}{#1}}
\newcommand{\CommentTok}[1]{\textcolor[rgb]{0.56,0.35,0.01}{\textit{#1}}}
\newcommand{\CommentVarTok}[1]{\textcolor[rgb]{0.56,0.35,0.01}{\textbf{\textit{#1}}}}
\newcommand{\ConstantTok}[1]{\textcolor[rgb]{0.56,0.35,0.01}{#1}}
\newcommand{\ControlFlowTok}[1]{\textcolor[rgb]{0.13,0.29,0.53}{\textbf{#1}}}
\newcommand{\DataTypeTok}[1]{\textcolor[rgb]{0.13,0.29,0.53}{#1}}
\newcommand{\DecValTok}[1]{\textcolor[rgb]{0.00,0.00,0.81}{#1}}
\newcommand{\DocumentationTok}[1]{\textcolor[rgb]{0.56,0.35,0.01}{\textbf{\textit{#1}}}}
\newcommand{\ErrorTok}[1]{\textcolor[rgb]{0.64,0.00,0.00}{\textbf{#1}}}
\newcommand{\ExtensionTok}[1]{#1}
\newcommand{\FloatTok}[1]{\textcolor[rgb]{0.00,0.00,0.81}{#1}}
\newcommand{\FunctionTok}[1]{\textcolor[rgb]{0.13,0.29,0.53}{\textbf{#1}}}
\newcommand{\ImportTok}[1]{#1}
\newcommand{\InformationTok}[1]{\textcolor[rgb]{0.56,0.35,0.01}{\textbf{\textit{#1}}}}
\newcommand{\KeywordTok}[1]{\textcolor[rgb]{0.13,0.29,0.53}{\textbf{#1}}}
\newcommand{\NormalTok}[1]{#1}
\newcommand{\OperatorTok}[1]{\textcolor[rgb]{0.81,0.36,0.00}{\textbf{#1}}}
\newcommand{\OtherTok}[1]{\textcolor[rgb]{0.56,0.35,0.01}{#1}}
\newcommand{\PreprocessorTok}[1]{\textcolor[rgb]{0.56,0.35,0.01}{\textit{#1}}}
\newcommand{\RegionMarkerTok}[1]{#1}
\newcommand{\SpecialCharTok}[1]{\textcolor[rgb]{0.81,0.36,0.00}{\textbf{#1}}}
\newcommand{\SpecialStringTok}[1]{\textcolor[rgb]{0.31,0.60,0.02}{#1}}
\newcommand{\StringTok}[1]{\textcolor[rgb]{0.31,0.60,0.02}{#1}}
\newcommand{\VariableTok}[1]{\textcolor[rgb]{0.00,0.00,0.00}{#1}}
\newcommand{\VerbatimStringTok}[1]{\textcolor[rgb]{0.31,0.60,0.02}{#1}}
\newcommand{\WarningTok}[1]{\textcolor[rgb]{0.56,0.35,0.01}{\textbf{\textit{#1}}}}
\usepackage{graphicx}
\makeatletter
\newsavebox\pandoc@box
\newcommand*\pandocbounded[1]{% scales image to fit in text height/width
  \sbox\pandoc@box{#1}%
  \Gscale@div\@tempa{\textheight}{\dimexpr\ht\pandoc@box+\dp\pandoc@box\relax}%
  \Gscale@div\@tempb{\linewidth}{\wd\pandoc@box}%
  \ifdim\@tempb\p@<\@tempa\p@\let\@tempa\@tempb\fi% select the smaller of both
  \ifdim\@tempa\p@<\p@\scalebox{\@tempa}{\usebox\pandoc@box}%
  \else\usebox{\pandoc@box}%
  \fi%
}
% Set default figure placement to htbp
\def\fps@figure{htbp}
\makeatother
\setlength{\emergencystretch}{3em} % prevent overfull lines
\providecommand{\tightlist}{%
  \setlength{\itemsep}{0pt}\setlength{\parskip}{0pt}}
\usepackage{bookmark}
\IfFileExists{xurl.sty}{\usepackage{xurl}}{} % add URL line breaks if available
\urlstyle{same}
\hypersetup{
  pdftitle={Eine Faktorenanalyse vom city lifestyle Datensatz},
  pdfauthor={Duc-Anh Nguyen},
  hidelinks,
  pdfcreator={LaTeX via pandoc}}

\title{Eine Faktorenanalyse vom city lifestyle Datensatz}
\author{Duc-Anh Nguyen}
\date{}

\begin{document}
\maketitle

\subsubsection{Ziel:}\label{ziel}

die zugrunde liegende Struktur der städtischen Lebensqualität: Wir
wollen verstehen, ob die Lebensqualität eher durch wirtschaftliche
Faktoren oder durch Umwelt-und Dichtefaktoren bestimmt wird. Damit
schaffen wir eine fundierte Datenbasis für weitere Surveys oder
politische Handlungsempfehlungen im Bereich der Stadtgeographie.''

\subsubsection{Datensatz}\label{datensatz}

Wir werden mit vielen verschiedenen Variablen arbeiten (Einkommen,
Miete, Luftqualität, Dichte, etc.). Den Datensatz kann man auf
\url{https://www.kaggle.com/datasets/umuttuygurr/city-lifestyle-segmentation-dataset}
finden.

\begin{Shaded}
\begin{Highlighting}[]
\NormalTok{df }\OtherTok{=} \FunctionTok{read.csv}\NormalTok{(}\StringTok{"../data/city\_lifestyle\_dataset.csv"}\NormalTok{)}
\FunctionTok{dim}\NormalTok{(df)}
\end{Highlighting}
\end{Shaded}

\begin{verbatim}
## [1] 300  10
\end{verbatim}

\begin{Shaded}
\begin{Highlighting}[]
\FunctionTok{head}\NormalTok{(df)}
\end{Highlighting}
\end{Shaded}

\begin{verbatim}
##       city_name country population_density avg_income internet_penetration
## 1     Old Vista  Europe               2775       3850                 86.4
## 2     Beachport  Europe               3861       3700                 78.1
## 3 Valleyborough  Europe               2562       4310                 80.1
## 4          City  Europe               3192       3970                 81.2
## 5         Falls  Europe               3496       4320                100.0
## 6   North Falls  Europe               1444       4640                 78.5
##   avg_rent air_quality_index public_transport_score happiness_score
## 1     1310                43                   52.0             8.5
## 2     1330                42                   62.8             8.1
## 3     1330                39                   73.2             8.5
## 4     1480                60                   49.2             8.5
## 5     1510                64                   93.7             8.5
## 6     1620                35                   51.6             8.4
##   green_space_ratio
## 1              23.8
## 2              33.1
## 3              40.2
## 4              43.6
## 5              42.5
## 6              20.3
\end{verbatim}

\begin{Shaded}
\begin{Highlighting}[]
\FunctionTok{colnames}\NormalTok{(df)}
\end{Highlighting}
\end{Shaded}

\begin{verbatim}
##  [1] "city_name"              "country"                "population_density"    
##  [4] "avg_income"             "internet_penetration"   "avg_rent"              
##  [7] "air_quality_index"      "public_transport_score" "happiness_score"       
## [10] "green_space_ratio"
\end{verbatim}

Wir nehmen an dass wir nur eine Stichprobe von N = 30 hätten, also
ziehen wie dann nur 30 Stichproben aus dem Datensatz zufällig aus. Dann
bearbeiten wir den Datensatz damit er einsatzbereits zum Analysieren
wäre.

\begin{Shaded}
\begin{Highlighting}[]
\FunctionTok{set.seed}\NormalTok{(}\DecValTok{123}\NormalTok{)}
\NormalTok{sample\_df }\OtherTok{\textless{}{-}}\NormalTok{  df }\SpecialCharTok{\%\textgreater{}\%} 
  \FunctionTok{sample\_n}\NormalTok{(}\DecValTok{30}\NormalTok{) }\SpecialCharTok{\%\textgreater{}\%} 
  \FunctionTok{select}\NormalTok{(}\StringTok{"population\_density"}\NormalTok{, }\StringTok{"avg\_income"}\NormalTok{, }\StringTok{"avg\_rent"}\NormalTok{, }\StringTok{"air\_quality\_index"}\NormalTok{, }\StringTok{"public\_transport\_score"}\NormalTok{, }\StringTok{"happiness\_score"}\NormalTok{, }\StringTok{"green\_space\_ratio"}\NormalTok{) }\SpecialCharTok{\%\textgreater{}\%} 
  \FunctionTok{rename}\NormalTok{(}
\NormalTok{    Bevölkerungsdichte }\OtherTok{=} \StringTok{"population\_density"}\NormalTok{, }
    \AttributeTok{Einkommen =} \StringTok{"avg\_income"}\NormalTok{, }
    \AttributeTok{Miete =} \StringTok{"avg\_rent"}\NormalTok{, }
    \AttributeTok{Luftverschmutzung =} \StringTok{"air\_quality\_index"}\NormalTok{, }
    \AttributeTok{OEPNV\_Score =} \StringTok{"public\_transport\_score"}\NormalTok{, }
    \AttributeTok{Zufriedenheit =} \StringTok{"happiness\_score"}\NormalTok{, }
    \AttributeTok{Gruenflaechen =} \StringTok{"green\_space\_ratio"}
\NormalTok{  )}

\CommentTok{\# Metadata}
\FunctionTok{attr}\NormalTok{(sample\_df, }\StringTok{"Beschreibung"}\NormalTok{) }\OtherTok{\textless{}{-}} \StringTok{"Pilotstudie zur Lebensqualität in 30 }
\StringTok{Städten"}
\FunctionTok{attr}\NormalTok{(sample\_df}\SpecialCharTok{$}\NormalTok{Bevölkerungsdichte, }\StringTok{"Beschreibung"}\NormalTok{) }\OtherTok{\textless{}{-}} \StringTok{"Einwohner pro Quadratkilometer (Bevölkerungsdichte)"}
\FunctionTok{attr}\NormalTok{(sample\_df}\SpecialCharTok{$}\NormalTok{Einkommen, }\StringTok{"Beschreibung"}\NormalTok{) }\OtherTok{\textless{}{-}} \StringTok{"Durchschnittliches Jahreseinkommen in USD"}
\FunctionTok{attr}\NormalTok{(sample\_df}\SpecialCharTok{$}\NormalTok{Miete, }\StringTok{"Beschreibung"}\NormalTok{) }\OtherTok{\textless{}{-}} \StringTok{"Durchschnittliche monatliche Mietkosten in USD"}
\FunctionTok{attr}\NormalTok{(sample\_df}\SpecialCharTok{$}\NormalTok{Luftverschmutzung, }\StringTok{"Beschreibung"}\NormalTok{) }\OtherTok{\textless{}{-}} \StringTok{"Air Quality Index (Höhere Werte = schlechtere Luft)"}
\FunctionTok{attr}\NormalTok{(sample\_df}\SpecialCharTok{$}\NormalTok{OEPNV\_Score, }\StringTok{"Beschreibung"}\NormalTok{) }\OtherTok{\textless{}{-}} \StringTok{"Index für die Qualität des öffentlichen Nahverkehrs (0{-}100)"}
\FunctionTok{attr}\NormalTok{(sample\_df}\SpecialCharTok{$}\NormalTok{Zufriedenheit, }\StringTok{"Beschreibung"}\NormalTok{) }\OtherTok{\textless{}{-}} \StringTok{"Selbsteingeschätzte Lebenszufriedenheit (Skala 1{-}10)"}
\FunctionTok{attr}\NormalTok{(sample\_df}\SpecialCharTok{$}\NormalTok{Gruenflaechen, }\StringTok{"Beschreibung"}\NormalTok{) }\OtherTok{\textless{}{-}} \StringTok{"Anteil der Grünflächen an der Gesamtfläche der Stadt (\%)"}
\FunctionTok{attr}\NormalTok{(sample\_df, }\StringTok{"codebook"}\NormalTok{) }\OtherTok{\textless{}{-}} \StringTok{"Daten stammen aus dem Kaggle City Lifestyle Dataset. N=30."}
\FunctionTok{attr}\NormalTok{(sample\_df, }\StringTok{"Kontakt"}\NormalTok{) }\OtherTok{\textless{}{-}} \StringTok{"anh.nguyen1@campus.lmu.de"}
\end{Highlighting}
\end{Shaded}

\subsubsection{Methodik:}\label{methodik}

Die Faktorenanalyse dient dazu, diese Komplexität zu reduzieren und die
latenten Dimensionen zu finden, die das Leben in einer Stadt tatsächlich
prägen:

\begin{enumerate}
\def\labelenumi{\arabic{enumi}.}
\item
  Komplexitätsreduktion (Variablen bündeln): Anstatt sieben einzelne
  Variablen getrennt zu betrachten, wollen wir sehen, welche Variablen
  zusammengehören. Wir wollen sehen, ob die Kovariablen stark
  korrelieren, dadurch später die Faktoren zu benennen, die man nicht
  direkt messen kann.
\item
  Statistische Validierung für kleine Stichproben (N=30) Mit einer
  kleinen Durch die Iteration und die Wahl der richtigen Rotation
  könnten wir zeigen, durch welche Faktoren der Unterschiede zwischen
  den Städten erklärt werden können. Wir führen zuvor einen Test bzgl.
  des Kaiser-Meyer-Olkin-Kriteriums, um zu sehen, ob wir überhaupt mit
  einer Hauptkomponentenanalyse fortfahren sollten. Das KMO-Kriterium
  wird aus den partiellen Korrelationen zwischen Itempaaren berechnet.
\end{enumerate}

Hier schauen wir uns erstmal die Korrelationen zwischen allen
Kovariablen.

\begin{Shaded}
\begin{Highlighting}[]
\CommentTok{\# scaled data for accurate Factor Analysis}
\NormalTok{sample\_df }\OtherTok{\textless{}{-}} \FunctionTok{scale}\NormalTok{(sample\_df)}

\NormalTok{cor\_mat }\OtherTok{\textless{}{-}} \FunctionTok{cor}\NormalTok{(sample\_df, }\AttributeTok{method =} \StringTok{"pearson"}\NormalTok{)}
\CommentTok{\# Bartlett{-}Test (Sind die Korrelationen signifikant von Null verschieden?)}
\NormalTok{bartlett\_test }\OtherTok{\textless{}{-}} \FunctionTok{cortest.bartlett}\NormalTok{(cor\_mat, }\AttributeTok{n =} \DecValTok{30}\NormalTok{)}
\FunctionTok{print}\NormalTok{(}\FunctionTok{paste}\NormalTok{(}\StringTok{"Bartlett{-}Test p{-}Wert:"}\NormalTok{, bartlett\_test}\SpecialCharTok{$}\NormalTok{p.value))}
\end{Highlighting}
\end{Shaded}

\begin{verbatim}
## [1] "Bartlett-Test p-Wert: 2.78121297788308e-28"
\end{verbatim}

\begin{Shaded}
\begin{Highlighting}[]
\FunctionTok{corrplot}\NormalTok{(cor\_mat, }\AttributeTok{method =} \StringTok{"color"}\NormalTok{, }\AttributeTok{addCoef.col =} \StringTok{"black"}\NormalTok{, }\AttributeTok{tl.col =} \StringTok{"black"}\NormalTok{, }\AttributeTok{title =} \StringTok{"Korrellation: Stichprobe (N = 30)"}\NormalTok{, }\AttributeTok{mar =} \FunctionTok{c}\NormalTok{(}\DecValTok{0}\NormalTok{,}\DecValTok{0}\NormalTok{,}\DecValTok{1}\NormalTok{,}\DecValTok{0}\NormalTok{))}
\end{Highlighting}
\end{Shaded}

\pandocbounded{\includegraphics[keepaspectratio]{main_files/figure-latex/correlation-1.pdf}}

\begin{Shaded}
\begin{Highlighting}[]
\CommentTok{\# Test mit Kaiser{-}Meyer{-}Olkin{-}Kriterium (Maß für die Eignung der Stichprobe)}
\CommentTok{\# Werte über 0.5 gelten als akzeptabel, über 0.7 als gut.}
\NormalTok{kmo\_test }\OtherTok{\textless{}{-}} \FunctionTok{KMO}\NormalTok{(sample\_df)}
\FunctionTok{print}\NormalTok{(kmo\_test}\SpecialCharTok{$}\NormalTok{MSA)}
\end{Highlighting}
\end{Shaded}

\begin{verbatim}
## [1] 0.574841
\end{verbatim}

Starke signifikante positive Korrelation zwischen Einkommen und Miete,
Einkommen und Zufriedenheit, Miete und Zufriedenheit, Luftverschmutzung
und Bevölkerungsdichte,

Starke signifikante negative Korrelation zwischen Luftverschmutzung und
Zufriedenheit, Luftverschmutzung und Grünflächen.

\subsubsection{PCA und Scree-Plot}\label{pca-und-scree-plot}

Wir führen hier vor der Faktorenanalyse eine PCA um herauszufinden, mit
vielen Komponenten wir arbeiten sollten.

\begin{Shaded}
\begin{Highlighting}[]
\CommentTok{\# PCA for Scree Plot}
\NormalTok{pca\_res }\OtherTok{\textless{}{-}} \FunctionTok{prcomp}\NormalTok{(sample\_df, }\AttributeTok{center =} \ConstantTok{TRUE}\NormalTok{, }\AttributeTok{scale. =} \ConstantTok{TRUE}\NormalTok{)}

\NormalTok{eigenvalues }\OtherTok{\textless{}{-}}\NormalTok{ pca\_res}\SpecialCharTok{$}\NormalTok{sdev}\SpecialCharTok{\^{}}\DecValTok{2}

\FunctionTok{plot}\NormalTok{(eigenvalues, }\AttributeTok{type =} \StringTok{"b"}\NormalTok{, }\AttributeTok{pch =} \FloatTok{0.19}\NormalTok{, }\AttributeTok{col =} \StringTok{"blue"}\NormalTok{, }
     \AttributeTok{main =} \StringTok{"Scree Plot: Stichprobe"}\NormalTok{, }
     \AttributeTok{xlab =} \StringTok{"Hauptkomponente"}\NormalTok{, }
     \AttributeTok{ylab =} \StringTok{"Eigenwert (Varianz)"}\NormalTok{) }
\FunctionTok{abline}\NormalTok{(}\AttributeTok{h =} \DecValTok{1}\NormalTok{, }\AttributeTok{col=}\StringTok{"red"}\NormalTok{, }\AttributeTok{lty=}\DecValTok{2}\NormalTok{, }\AttributeTok{lwd =} \DecValTok{2}\NormalTok{) }\CommentTok{\# Kaiser Criterion}
\FunctionTok{legend}\NormalTok{(}\StringTok{"topright"}\NormalTok{, }\AttributeTok{legend =} \StringTok{"Kaiser{-}Kriterium (EV = 1)"}\NormalTok{, }\AttributeTok{col =} \StringTok{"red"}\NormalTok{, }\AttributeTok{lty =} \DecValTok{2}\NormalTok{, }\AttributeTok{lwd =} \DecValTok{2}\NormalTok{)}
\end{Highlighting}
\end{Shaded}

\pandocbounded{\includegraphics[keepaspectratio]{main_files/figure-latex/pca-1.pdf}}

Bei dem Elbogen hätten wir 2 Hauptkomponente.

\subsubsection{Final Factor Analysis}\label{final-factor-analysis}

Eine Faktorenanalyse mit Varimax Rotation und dem
Minimum-Residual-Ansatz (minres) um latenten Strukturen zu
identifizieren.

Wir testen nun iterativ ein 2-Faktoren-Modell gegen ein
3-Faktoren-Modell und messen die Signifikanz anhand des RMSR (Root Mean
Square Residual) und des p-Wertes des Modells.

\begin{Shaded}
\begin{Highlighting}[]
\CommentTok{\# Iteration 1: Das 2{-}Faktoren{-}Modell (MinRes)}
\NormalTok{fa\_2 }\OtherTok{\textless{}{-}} \FunctionTok{fa}\NormalTok{(sample\_df, }\AttributeTok{nfactors =} \DecValTok{2}\NormalTok{, }\AttributeTok{rotate =} \StringTok{"varimax"}\NormalTok{, }\AttributeTok{fm =} \StringTok{"minres"}\NormalTok{)}

\CommentTok{\# Iteration 2: Das 3{-}Faktoren{-}Modell (MinRes)}
\NormalTok{fa\_3 }\OtherTok{\textless{}{-}} \FunctionTok{fa}\NormalTok{(sample\_df, }\AttributeTok{nfactors =} \DecValTok{3}\NormalTok{, }\AttributeTok{rotate =} \StringTok{"varimax"}\NormalTok{, }\AttributeTok{fm =} \StringTok{"minres"}\NormalTok{)}

\CommentTok{\# Signifikanz und Modellgüte (Fit) vergleichen:}
\FunctionTok{cat}\NormalTok{(}\StringTok{"}\SpecialCharTok{\textbackslash{}n}\StringTok{{-}{-}{-} MODELL{-}VERGLEICH {-}{-}{-}}\SpecialCharTok{\textbackslash{}n}\StringTok{"}\NormalTok{)}
\end{Highlighting}
\end{Shaded}

\begin{verbatim}
## 
## --- MODELL-VERGLEICH ---
\end{verbatim}

\begin{Shaded}
\begin{Highlighting}[]
\FunctionTok{cat}\NormalTok{(}\StringTok{"2 Faktoren RMSR (kleiner ist besser):"}\NormalTok{, }\FunctionTok{round}\NormalTok{(fa\_2}\SpecialCharTok{$}\NormalTok{rms, }\DecValTok{3}\NormalTok{), }\StringTok{"}\SpecialCharTok{\textbackslash{}n}\StringTok{"}\NormalTok{)}
\end{Highlighting}
\end{Shaded}

\begin{verbatim}
## 2 Faktoren RMSR (kleiner ist besser): 0.05
\end{verbatim}

\begin{Shaded}
\begin{Highlighting}[]
\FunctionTok{cat}\NormalTok{(}\StringTok{"3 Faktoren RMSR (kleiner ist besser):"}\NormalTok{, }\FunctionTok{round}\NormalTok{(fa\_3}\SpecialCharTok{$}\NormalTok{rms, }\DecValTok{3}\NormalTok{), }\StringTok{"}\SpecialCharTok{\textbackslash{}n}\StringTok{"}\NormalTok{)}
\end{Highlighting}
\end{Shaded}

\begin{verbatim}
## 3 Faktoren RMSR (kleiner ist besser): 0.02
\end{verbatim}

\begin{Shaded}
\begin{Highlighting}[]
\FunctionTok{cat}\NormalTok{(}\StringTok{"2 Faktoren Tucker{-}Lewis Index (TLI, \textgreater{} 0.9 ist gut):"}\NormalTok{, }\FunctionTok{round}\NormalTok{(fa\_2}\SpecialCharTok{$}\NormalTok{TLI, }\DecValTok{3}\NormalTok{), }\StringTok{"}\SpecialCharTok{\textbackslash{}n}\StringTok{"}\NormalTok{)}
\end{Highlighting}
\end{Shaded}

\begin{verbatim}
## 2 Faktoren Tucker-Lewis Index (TLI, > 0.9 ist gut): 0.505
\end{verbatim}

\begin{Shaded}
\begin{Highlighting}[]
\CommentTok{\# Statistische Signifikanz des 2{-}Faktor{-}Modells (p{-}Wert)}
\CommentTok{\# Testet die Nullhypothese, dass 2 Faktoren ausreichen. (p \textgreater{} 0.05 ist hier GUT!)}
\FunctionTok{cat}\NormalTok{(}\StringTok{"Signifikanz des 2{-}Faktor{-}Modells (p{-}Wert):"}\NormalTok{, fa\_2}\SpecialCharTok{$}\NormalTok{PVAL, }\StringTok{"}\SpecialCharTok{\textbackslash{}n}\StringTok{"}\NormalTok{)}
\end{Highlighting}
\end{Shaded}

\begin{verbatim}
## Signifikanz des 2-Faktor-Modells (p-Wert): 1.065288e-05
\end{verbatim}

RMSR (0.05 vs.~0.02): Der Root Mean Square Residual misst den
durchschnittlichen Fehler: der Wert von 0.05 beim 2-Faktoren-Modell ist
gut. Dass er beim 3-Faktoren-Modell auf 0.02 sinkt, ist logisch: Je mehr
Faktoren man hinzufügt, desto besser schmiegt sich das Modell an die
Daten an.

TLI (0.505): Der Tucker-Lewis Index hat den Wert von 0.505 ist sehr
niedrig. Das bedeutet: das Modell erklärt die Daten deutlich besser als
ein Zufallsmodell, aber es gibt noch viel ungeklärte Struktur.

Signifikanz (p-Wert = 0): In der Faktorenanalyse ist ein p-Wert von
1.065288e-05 (signifikant) paradoxerweise oft ein Zeichen, dass das
Modell nicht perfekt passt. Wir lehnen die Nullhypothese ab, dass das
Modell die Daten exakt abbildet.

\subsubsection{Interpretation:}\label{interpretation}

\begin{Shaded}
\begin{Highlighting}[]
\FunctionTok{print}\NormalTok{(fa\_2}\SpecialCharTok{$}\NormalTok{loadings, }\AttributeTok{cutoff =} \FloatTok{0.3}\NormalTok{)}
\end{Highlighting}
\end{Shaded}

\begin{verbatim}
## 
## Loadings:
##                    MR1    MR2   
## Bevölkerungsdichte         0.787
## Einkommen           0.940       
## Miete               0.900       
## Luftverschmutzung  -0.428  0.801
## OEPNV_Score         0.629  0.322
## Zufriedenheit       0.835 -0.471
## Gruenflaechen             -0.750
## 
##                  MR1   MR2
## SS loadings    2.970 2.200
## Proportion Var 0.424 0.314
## Cumulative Var 0.424 0.739
\end{verbatim}

Faktor/MR 1: Sozioökonomische Lebensqualität: Key Loadings: Einkommen
(0.94), Miete(0.9), Zufriedenheit(0.84) und OEPNV\_Score(0.63). Dieses
Faktoren decken 42\% der Varianz ab, und könnten sich so interpretieren
lassen: Städte mit höheren Einkommen und Miete hätten auch ``besseren''
OEPNV und zufriedeneren BewöhnerInnen. Mit -0.43 bei Luftverschmutzung
suggeriert, dass solche Städte auch wenig verschmutzt sind.

Faktor/MR 2: Ubarner Belastungsfaktor: Key Loadings: Luftverschmutzung
(0.8), Bevölkerungsdichte(0.79), Gruenflaechen (-0.75) und Zufriedenheit
(-0.47). Diese Faktoren repräsentieren die Kosten einer Stadt: Wir
sehen, dass höhere Bevölkerungsdichte kommt häufig mit höherer
Luftverschmutzung und weinigeren Gruenflaechen, welche die Zufriedenheit
der EinwöhnerInnen senken.

Das Modell ist robust: wir erreichen eine kumulierte Varainzaufklärung
von fast 74\%, die zeigt, dass die 2 Faktoren die wesentlichen
Unterschiede in der Lebensqualität unserer 30 Städte sehr gut erfassen.
Ein interessanter Punkt ist die Luftverschmutzung: Sie lädt negativ auf
dem ersten Faktor (-0.43) und stark positiv auf dem zweiten (0.80). Das
zeigt, dass Umweltbelastung in meinem Modell zwei Seiten hat: Sie sinkt
mit steigendem Wohlstand, steigt aber massiv mit der baulichen Dichte.

\end{document}
